% ============================================================
\chapter{Introduction au TAL}
\label{ch:introduction}
% ============================================================

En 1950, Alan Turing pose la question fondatrice de l'intelligence
artificielle~: une machine peut-elle \og penser\fg{}~? Et il propose un
crit\`ere op\'erationnel~: si un humain, en conversant avec une machine par
texte interpos\'e, ne parvient pas \`a la distinguer d'un autre humain, alors
la machine \og pense\fg{} au sens fonctionnel du terme. Ce \og test de
Turing\fg{} est fondamentalement un d\'efi de \emph{traitement du langage
naturel}~: pour le r\'eussir, une machine doit comprendre la syntaxe, la
s\'emantique, le contexte, l'ironie, le sous-entendu. Soixante-quinze ans
plus tard, les mod\`eles de langage de grande taille (GPT, Claude, LLaMA)
semblent s'en approcher --- mais les d\'efis fondamentaux restent
intacts.

\begin{intuition}[title=\textbf{Pourquoi \'etudier le TAL ?}]
Le langage est le principal vecteur de communication humaine. Enseigner aux
machines \`a comprendre, produire et raisonner sur le langage naturel est l'un
des d\'efis fondamentaux de l'intelligence artificielle. Les applications vont
de la traduction automatique \`a l'extraction d'information, en passant par les
agents conversationnels et la g\'en\'eration de texte.
\end{intuition}

% ------------------------------------------------------------
\section{Qu'est-ce que le TAL ?}
\label{sec:intro:quoi}
% ------------------------------------------------------------

\begin{definition}[Traitement Automatique du Langage]
Le \emph{Traitement Automatique du Langage} (TAL) est le sous-domaine de
l'intelligence artificielle qui \'etudie les interactions entre les ordinateurs
et le langage humain. Il englobe l'analyse, la compr\'ehension et la g\'en\'eration
de texte et de parole en langues naturelles.
\end{definition}

Le TAL se distingue de la \emph{linguistique computationnelle} par son
orientation vers les applications, bien que les deux domaines partagent un
socle th\'eorique commun. Formellement, on peut consid\'erer le langage comme
un syst\`eme g\'en\'eratif d\'efini sur un vocabulaire $\mathcal{V}$ :
\[
\mathcal{L} \subseteq \mathcal{V}^* = \bigcup_{n=0}^{\infty} \mathcal{V}^n
\]
o\`u $\mathcal{V}^*$ d\'esigne l'ensemble de toutes les s\'equences finies
sur $\mathcal{V}$.

% ------------------------------------------------------------
\section{Niveaux d'analyse linguistique}
\label{sec:intro:niveaux}
% ------------------------------------------------------------

Le traitement du langage naturel op\`ere \`a plusieurs niveaux d'analyse :

\begin{enumerate}
    \item \textbf{Phon\'etique et phonologie} : \'etude des sons du langage
          (principalement pour le traitement de la parole).
    \item \textbf{Morphologie} : structure interne des mots (pr\'efixes,
          suffixes, flexions). Par exemple, \og anti-constitutionnel\fg
          se d\'ecompose en \texttt{anti-} + \texttt{constitution} +
          \texttt{-nel}.
    \item \textbf{Syntaxe} : structure des phrases, relations grammaticales.
    \item \textbf{S\'emantique} : sens des mots et des phrases.
    \item \textbf{Pragmatique} : sens en contexte, intentions du locuteur.
    \item \textbf{Discours} : coh\'erence et structure au-del\`a de la phrase.
\end{enumerate}

\begin{remarque}
En TAL moderne, les mod\`eles profonds apprennent implicitement ces niveaux
d'analyse \`a partir des donn\'ees, sans segmentation explicite. Cependant,
comprendre ces niveaux reste essentiel pour diagnostiquer les erreurs et
concevoir des \'evaluations pertinentes.
\end{remarque}

% ------------------------------------------------------------
\section{T\^aches fondamentales du TAL}
\label{sec:intro:taches}
% ------------------------------------------------------------

\subsection{Classification de texte}

La classification de texte consiste \`a assigner une \'etiquette $y \in \mathcal{Y}$
\`a un document $\mathbf{x} = (x_1, x_2, \ldots, x_T)$. Le mod\`ele apprend
une fonction $f : \mathcal{V}^* \to \mathcal{Y}$ \`a partir d'un ensemble
d'entra\^inement $\{(\mathbf{x}_i, y_i)\}_{i=1}^{N}$.

Applications : analyse de sentiments, d\'etection de spam, classification
th\'ematique.

\subsection{\'Etiquetage de s\'equences}

L'\'etiquetage de s\'equences assigne une \'etiquette $y_t \in \mathcal{Y}$
\`a chaque token $x_t$ d'une s\'equence :
\[
f : \mathcal{V}^T \to \mathcal{Y}^T
\]

Applications : reconnaissance d'entit\'es nomm\'ees (NER), \'etiquetage
morpho-syntaxique (POS tagging).

\begin{exemple}
Consid\'erons la phrase \og Marie habite \`a Paris \fg avec l'\'etiquetage
NER suivant :
\begin{center}
\begin{tabular}{ccccc}
Marie & habite & \`a & Paris \\
\texttt{B-PER} & \texttt{O} & \texttt{O} & \texttt{B-LOC}
\end{tabular}
\end{center}
\end{exemple}

\subsection{Traduction automatique}

La traduction automatique est une t\^ache de transduction s\'equence-\`a-s\'equence :
\[
f : \mathcal{V}_{\text{source}}^* \to \mathcal{V}_{\text{cible}}^*
\]
C'est historiquement l'une des t\^aches motrices du TAL.

\subsection{Question-r\'eponse}

Soit un contexte $C$ et une question $Q$, le syst\`eme doit extraire ou
g\'en\'erer une r\'eponse $A$ :
\[
f : (C, Q) \to A
\]

\subsection{R\'esum\'e automatique}

Le r\'esum\'e automatique produit une version condens\'ee $S$ d'un document $D$,
pr\'eservant l'information essentielle : $f : D \to S$ avec $\abs{S} \ll \abs{D}$.

% ------------------------------------------------------------
\section{Fondements probabilistes}
\label{sec:intro:proba}
% ------------------------------------------------------------

La mod\'elisation probabiliste du langage est au c\oe ur du TAL. Un
\emph{mod\`ele de langue} d\'efinit une distribution de probabilit\'e sur les
s\'equences de tokens.

\begin{definition}[Mod\`ele de langue]
Un \emph{mod\`ele de langue} est une distribution de probabilit\'e $P$ sur
$\mathcal{V}^*$ telle que :
\[
P(\mathbf{x}) = P(x_1, x_2, \ldots, x_T) = \prod_{t=1}^{T} P(x_t \mid x_1, \ldots, x_{t-1})
\]
o\`u la factorisation d\'ecoule de la r\`egle de la cha\^ine des probabilit\'es.
\end{definition}

\begin{theoreme}[R\`egle de la cha\^ine]
Pour toute distribution jointe $P(x_1, \ldots, x_T)$ :
\[
P(x_1, \ldots, x_T) = \prod_{t=1}^{T} P(x_t \mid x_{<t})
\]
o\`u $x_{<t} = (x_1, \ldots, x_{t-1})$ d\'esigne le \emph{contexte} au pas $t$.
\end{theoreme}

\begin{proof}
Par applications successives de la d\'efinition de la probabilit\'e conditionnelle
$P(A \mid B) = P(A, B) / P(B)$ :
\begin{align*}
P(x_1, x_2, \ldots, x_T) &= P(x_1) \cdot P(x_2 \mid x_1) \cdot P(x_3 \mid x_1, x_2) \cdots \\
                           &= \prod_{t=1}^{T} P(x_t \mid x_{<t}). \qedhere
\end{align*}
\end{proof}

% ------------------------------------------------------------
\section{Th\'eorie de l'information et langage}
\label{sec:intro:info}
% ------------------------------------------------------------

\begin{definition}[Entropie]
L'\emph{entropie} d'une variable al\'eatoire discr\`ete $X$ \`a valeurs dans
$\mathcal{X}$ est :
\[
H(X) = -\sum_{x \in \mathcal{X}} P(x) \log_2 P(x)
\]
Elle mesure l'incertitude moyenne associ\'ee \`a $X$.
\end{definition}

\begin{definition}[Entropie crois\'ee]
L'\emph{entropie crois\'ee} entre la distribution r\'eelle $p$ et un mod\`ele $q$ est :
\[
H(p, q) = -\sum_{x} p(x) \log q(x) = H(p) + \KL(p \| q)
\]
\end{definition}

\begin{proposition}[Borne inf\'erieure]
Pour tout mod\`ele $q$, on a $H(p, q) \geq H(p)$, avec \'egalit\'e si et
seulement si $q = p$.
\end{proposition}

La \emph{perplexit\'e} d'un mod\`ele de langue $q$ est d\'efinie par :
\[
\text{PPL}(q) = 2^{H(p, q)}
\]

\begin{formules}[title=\textbf{M\'etriques informationnelles cl\'es}]
\begin{align}
H(X) &= -\sum_{x} P(x) \log P(x) & \text{(entropie)} \\
\KL(p \| q) &= \sum_{x} p(x) \log \frac{p(x)}{q(x)} & \text{(divergence KL)} \\
\text{PPL} &= \exp\!\left(-\frac{1}{T}\sum_{t=1}^{T} \log q(x_t \mid x_{<t})\right) & \text{(perplexit\'e)}
\end{align}
\end{formules}

% ------------------------------------------------------------
\section{Pipeline classique de TAL}
\label{sec:intro:pipeline}
% ------------------------------------------------------------

Un pipeline de TAL classique comprend les \'etapes suivantes :

\begin{enumerate}
    \item \textbf{Collecte de donn\'ees} : corpus textuels, web scraping, APIs.
    \item \textbf{Pr\'etraitement} :
    \begin{itemize}
        \item Tokenisation (d\'ecoupage en unit\'es lexicales)
        \item Normalisation (mise en minuscules, suppression d'accents)
        \item Suppression des mots vides (\emph{stopwords})
        \item Racinisation (\emph{stemming}) ou lemmatisation
    \end{itemize}
    \item \textbf{Repr\'esentation} : conversion du texte en vecteurs num\'eriques.
    \item \textbf{Mod\'elisation} : entra\^inement d'un mod\`ele.
    \item \textbf{\'Evaluation} : mesure de la performance sur un jeu de test.
    \item \textbf{D\'eploiement} : mise en production du mod\`ele.
\end{enumerate}

\begin{attention}[title=\textbf{Pr\'etraitement et mod\`eles modernes}]
Avec les mod\`eles pr\'e-entra\^in\'es (BERT, GPT), le pr\'etraitement est
g\'en\'eralement minimal : le tokeniseur du mod\`ele (BPE, WordPiece) remplace
les \'etapes classiques de stemming et de suppression de stopwords. Appliquer
un pr\'etraitement agressif peut m\^eme d\'egrader les performances.
\end{attention}

% ------------------------------------------------------------
\section{Tokenisation}
\label{sec:intro:tokenisation}
% ------------------------------------------------------------

\begin{definition}[Tokenisation]
La \emph{tokenisation} est le processus de d\'ecoupage d'un texte en unit\'es
\'el\'ementaires appel\'ees \emph{tokens}. Un token peut \^etre un mot, un
sous-mot, un caract\`ere ou un byte.
\end{definition}

\subsection{Tokenisation par sous-mots}

Les algorithmes modernes de tokenisation par sous-mots (BPE, WordPiece,
Unigram) offrent un compromis entre la tokenisation par mots (vocabulaire
tr\`es grand) et par caract\`eres (s\'equences tr\`es longues).

\begin{algorithme}[title=\textbf{Byte Pair Encoding (BPE)}]
\begin{enumerate}
    \item Initialiser le vocabulaire avec tous les caract\`eres individuels.
    \item R\'ep\'eter $k$ fois :
    \begin{enumerate}
        \item Compter toutes les paires de symboles adjacents dans le corpus.
        \item Fusionner la paire la plus fr\'equente en un nouveau symbole.
        \item Ajouter ce symbole au vocabulaire.
    \end{enumerate}
    \item Retourner le vocabulaire final de taille $\abs{\mathcal{V}_0} + k$.
\end{enumerate}
\end{algorithme}

\begin{codeblock}[title=\textbf{Tokenisation avec Hugging Face}]
\begin{minted}{python}
from transformers import AutoTokenizer

tokenizer = AutoTokenizer.from_pretrained("bert-base-uncased")
text = "Natural language processing is fascinating."
tokens = tokenizer.tokenize(text)
print(tokens)
# ['natural', 'language', 'processing', 'is', 'fascinating', '.']

ids = tokenizer.encode(text, return_tensors="pt")
print(ids)
# tensor([[ 101, 3019, 2653, 6364, 2003, 15281, 1012,  102]])
\end{minted}
\end{codeblock}

% ------------------------------------------------------------
\section{Bref historique du TAL}
\label{sec:intro:histoire}
% ------------------------------------------------------------

\begin{center}
\renewcommand{\arraystretch}{1.2}
\begin{tabular}{lll}
\toprule
\textbf{P\'eriode} & \textbf{Paradigme} & \textbf{Exemples} \\
\midrule
1950--1970 & R\`egles symboliques & ELIZA, SHRDLU \\
1970--1990 & Grammaires formelles & ATN, LFG, HPSG \\
1990--2010 & M\'ethodes statistiques & HMM, CRF, $n$-grammes \\
2010--2017 & Apprentissage profond & Word2Vec, LSTM, Seq2Seq \\
2017--      & Transformers \& LLM & BERT, GPT, T5, LLaMA \\
\bottomrule
\end{tabular}
\end{center}

\begin{remarque}
Le passage du paradigme statistique au paradigme neuronal ne repr\'esente pas
une rupture totale : les concepts d'entropie crois\'ee, de mod\`ele de langue
et de r\'egularisation restent au c\oe ur des approches modernes.
\end{remarque}

% ------------------------------------------------------------
\section{D\'efis ouverts}
\label{sec:intro:defis}
% ------------------------------------------------------------

Malgr\'e les progr\`es spectaculaires, de nombreux d\'efis restent ouverts :

\begin{itemize}
    \item \textbf{Compr\'ehension r\'eelle vs.\ corr\'elations statistiques} :
          les LLM capturent-ils le \emph{sens} ou simplement des
          r\'egularit\'es de surface ?
    \item \textbf{Langues peu dott\'ees} : la majorit\'e des ressources sont
          disponibles en anglais ; environ 7\,000 langues sont parl\'ees dans
          le monde.
    \item \textbf{Robustesse} : sensibilit\'e aux perturbations adversariales,
          aux n\'eologismes, aux domaines hors distribution.
    \item \textbf{Efficacit\'e} : les LLM requi\`erent des ressources
          computationnelles consid\'erables.
    \item \textbf{\'Ethique} : biais, toxicit\'e, confidentialit\'e,
          d\'esinformation.
\end{itemize}

% ------------------------------------------------------------
\section{Exercices}
\label{sec:intro:exercices}
% ------------------------------------------------------------

\begin{exercice}[$\star$]
Calculez l'entropie d'un mod\`ele de langue uniforme sur un vocabulaire de
taille $V = 50\,000$. Quelle est la perplexit\'e correspondante ?
\end{exercice}

\begin{exercice}[$\star$]
Soit la phrase \og Le chat mange la souris \fg. Appliquez la r\`egle de
la cha\^ine pour exprimer $P(\text{Le, chat, mange, la, souris})$ comme un
produit de probabilit\'es conditionnelles.
\end{exercice}

\begin{exercice}[$\star\star$]
Impl\'ementez un tokeniseur BPE simple en Python. Partez d'un petit corpus
de 10 phrases et montrez l'\'evolution du vocabulaire apr\`es 20 fusions.
\end{exercice}

\begin{exercice}[$\star\star$]
Montrez que la divergence de Kullback-Leibler $\KL(p \| q)$ est toujours
positive ou nulle (in\'egalit\'e de Gibbs). \emph{Indication :} utilisez
l'in\'egalit\'e $\log x \leq x - 1$.
\end{exercice}

\begin{exercice}[$\star\star\star$]
Comparez la perplexit\'e d'un mod\`ele $n$-gramme (pour $n = 1, 2, 3$) et
d'un mod\`ele neuronal (GPT-2) sur un corpus de votre choix. Analysez les
r\'esultats en fonction de la taille du corpus d'entra\^inement.
\end{exercice}
